%% exemplo-guia-visual.tex
%% Documento mínimo de teste — linha Guia Visual 16:9 CRP-SP
%% Compilar 2x com: lualatex exemplo-guia-visual.tex
%% (a primeira passagem grava o .nav; a segunda renderiza a navbar)

\DocumentMetadata{
  lang        = pt-BR,
  pdfversion  = 2.0,
  pdfstandard = ua-2,
  testphase   = {phase-III},
}

\documentclass{crpsp-guia_visual}

% --- Paleta desta publicação ---
\crpspCorPrincipal{004A8F}
\crpspCorSecundaria{E8452A}

\hypersetup{
  pdftitle  = {Exemplo de Guia Visual CRP-SP},
  pdfauthor = {Conselho Regional de Psicologia de São Paulo},
}

\begin{document}

% ══════════════════════════════════════════════
\section[Apresentação]{Apresentação institucional do CRP-SP}
% ══════════════════════════════════════════════

\begin{frame}
  \begin{colunas}
    \coluna[span=2]{%
      Conteúdo principal do primeiro frame — ocupa dois terços da
      largura útil. O texto flui normalmente dentro da coluna, com
      a tipografia padrão da linha editorial.

      \begin{itemize}
        \item Primeiro item com bullet na cor principal.
        \item Segundo item da lista.
      \end{itemize}
    }%
    \coluna[span=1]{%
      % Placeholder de imagem: retângulo cheio (sem arquivo externo).
      % Em produção: \includegraphics[width=\linewidth, alt={...}]{imagem.pdf}
      \textcolor{cor3}{\rule{\linewidth}{45mm}}%
    }%
  \end{colunas}
\end{frame}

\begin{frame}[layout=full]
  Frame de largura total, sem grid de colunas. Serve para imagens
  panorâmicas ou blocos de texto únicos.

  \textcolor{cor2}{\rule{\linewidth}{30mm}}
\end{frame}

% ══════════════════════════════════════════════
\section[Resultados]{Resultados e indicadores de 2025}
% ══════════════════════════════════════════════

\begin{frame}
  \begin{colunas}
    \coluna[span=1]{%
      Primeira coluna do grid 1+1+1.

      \textcolor{cor1}{\rule{\linewidth}{25mm}}%
    }%
    \coluna[span=1]{%
      Segunda coluna, mesma largura.

      \textcolor{cor2}{\rule{\linewidth}{25mm}}%
    }%
    \coluna[span=1]{%
      Terceira coluna, fechando a soma de spans em 3.

      \textcolor{cor3}{\rule{\linewidth}{25mm}}%
    }%
  \end{colunas}
\end{frame}

\begin{frame}
  \begin{colunas}
    \coluna[span=1]{%
      Coluna estreita com legenda ou destaque lateral.
    }%
    \coluna[span=2]{%
      Coluna larga com o conteúdo principal deste frame — grid 1+2,
      espelhando o 2+1 do primeiro frame.
    }%
  \end{colunas}
\end{frame}

% Teste das chaves de alinhamento: as três colunas têm o mesmo texto
% curto e devem aparecer no topo, no meio e no pé da mancha.
\begin{frame}
  \begin{colunas}
    \coluna[span=1, t]{%
      Coluna \texttt{t}: conteúdo no topo (padrão).
    }%
    \coluna[span=1, c, centro]{%
      Coluna \texttt{c} com \texttt{centro}: centrada nos dois eixos.

      \textcolor{cor2}{\rule{20mm}{15mm}}%
    }%
    \coluna[span=1, b]{%
      Coluna \texttt{b}: conteúdo rente ao rodapé da mancha.
    }%
  \end{colunas}
\end{frame}

% ══════════════════════════════════════════════
\section{Encaminhamentos e próximos passos}
% (sem argumento curto: o título longo vai à navbar — teste de fallback)
% ══════════════════════════════════════════════

\begin{frame}[layout=full]
  Última seção, declarada sem título curto para testar o fallback
  da navbar.
\end{frame}

\end{document}
